% !TeX root = ../../Book1.tex
\section{\texorpdfstring{\(\pi\)}{π}-系统、\texorpdfstring{\(\lambda\)}{λ}-系统与单调类}
\label{b1:sec:0103}

生成定义把一个庞大的 \(\sigma\)-代数压缩到少量生成元，但还留下一个实际困难：为了证明某种性质对 \(\sigma(\mathcal E)\) 中所有集合成立，若把满足该性质的集合记为 \(\mathcal D\)，直接验证 \(\mathcal D\) 是 \(\sigma\)-代数往往要求过强。更常见的情形是，\(\mathcal D\) 只对差与不交并封闭，或只对单调极限封闭。本节的两个推广定理说明，只要生成元本身具有适当的有限运算结构，这些较弱闭性已经足够。

\subsection{\texorpdfstring{\(\pi\)}{π}-系统}
\label{b1:subsec:010301}

许多自然生成族不对补或并封闭，却对有限交封闭。区间相交仍是同类区间或空集，矩形相交仍是同类矩形或空集；这一条简单闭性正是后面论证所需的有限结构。

\begin{definition}
\label{b1:def:pi-system}
非空集合族 \(\mathcal C\subseteq\mathcal P(X)\) 若对任意 \(A,B\in\mathcal C\) 都有 \(A\cap B\in\mathcal C\)，则称为\term[pi-xitong]{\(\pi\)-系统}[pi-system]。
\end{definition}

例如，实直线上的闭半直线族
\(
\{(-\infty,t]:t\in\mathbb R\}
\)
是 \(\pi\)-系统，因为两个成员的交取较小端点。若 \(\mathcal A\) 与 \(\mathcal B\) 分别是 \(X\) 与 \(Y\) 上的集合代数，则矩形族
\[
\{\,A\times B:A\in\mathcal A,\ B\in\mathcal B\,\}
\]
也是 \(\pi\)-系统，这是因为
\[
(A_1\times B_1)\cap(A_2\times B_2)
=(A_1\cap A_2)\times(B_1\cap B_2).
\]
另一方面，所有开区间组成的集合族若不把空集计入，就不是 \(\pi\)-系统：两个不相交开区间的交为空集。使用 \(\pi\)-系统时不能只凭“形状相交后不变”而忽略退化情形。

\subsection{Dynkin \texorpdfstring{\(\lambda\)}{λ}-系统}
\label{b1:subsec:010302}

\(\sigma\)-代数要求任意可数并封闭，而很多等式只能自然地对两两不交的集合求和。Dynkin \(\lambda\)-系统把公理调整为“嵌套差与不交可数并”，因而特别适合承载这类性质。

\begin{definition}
\label{b1:def:lambda-system}
\term[dynkin-lambda-xitong]{Dynkin \(\lambda\)-系统}[Dynkin lambda-system]是集合族 \(\mathcal D\subseteq\mathcal P(X)\)，满足：
\begin{enumerate}
  \item \label{b1:item:lambda-system-1}
  \(X\in\mathcal D\)；
  \item \label{b1:item:lambda-system-2}
  若 \(A,B\in\mathcal D\) 且 \(A\subseteq B\)，则 \(B\setminus A\in\mathcal D\)；
  \item \label{b1:item:lambda-system-3}
  若 \(A_1,A_2,\ldots\in\mathcal D\) 两两不交，则 \(\bigcup_{n=1}^{\infty}A_n\in\mathcal D\)。
\end{enumerate}
\end{definition}

对任意集合族 \(\mathcal E\subseteq\mathcal P(X)\)，所有包含 \(\mathcal E\) 的 \(\lambda\)-系统之交仍是 \(\lambda\)-系统。把这个交记为 \(\lambda(\mathcal E)\)，称为由 \(\mathcal E\) 生成的 \(\lambda\)-系统。幂集保证参与求交的集合族非空。

\begin{proposition}
\label{b1:prop:lambda-basic-closure}
若 \(\mathcal D\) 是 \(\lambda\)-系统，则 \(\varnothing\in\mathcal D\)，并且 \(\mathcal D\) 对补集、递增并和递减交封闭。
\end{proposition}
\begin{proof}
在定义的第二条中取 \(A=B=X\)，得到 \(\varnothing\in\mathcal D\)；取 \(B=X\)，得到 \(X\setminus A=A^{\mathrm c}\in\mathcal D\)。

若 \(A_n\uparrow A\)，令
\[
B_1=A_1,
\quad
B_n=A_n\setminus A_{n-1}\quad(n\geq2).
\]
由嵌套差的封闭性，每个 \(B_n\in\mathcal D\)；这些集合两两不交，且 \(\bigcup_nB_n=\bigcup_nA_n=A\)。所以 \(A\in\mathcal D\)。若 \(A_n\downarrow A\)，则 \(A_n^{\mathrm c}\uparrow A^{\mathrm c}\)，先使用补集与递增并的封闭性得到 \(A^{\mathrm c}\in\mathcal D\)，再取一次补便有 \(A\in\mathcal D\)。
\end{proof}

每个 \(\sigma\)-代数都是 \(\lambda\)-系统，反过来却不成立。令
\[
X=\{1,2,3,4\},
\quad
\mathcal D
=\{\varnothing,X,\{1,2\},\{3,4\},\{1,3\},\{2,4\}\}.
\]
除相等情形外，\(\mathcal D\) 中两个嵌套成员只能以下端为空集或以上端为 \(X\)，故其差仍在 \(\mathcal D\) 中。两个非空、非全集的成员只有在互补时才不交，此时其并为 \(X\)。所以 \(\mathcal D\) 是 \(\lambda\)-系统。然而
\(
\{1,2\}\cap\{1,3\}=\{1\}\notin\mathcal D
\)，所以它不是 \(\pi\)-系统，更不是 \(\sigma\)-代数。缺失的有限交正是下一定理需要由生成元补回的结构。

\subsection{\texorpdfstring{\(\pi\)}{π}-\texorpdfstring{\(\lambda\)}{λ} 定理}
\label{b1:subsec:010303}

\begin{theorem}
\label{b1:thm:pi-lambda}
若 \(\mathcal P\) 是 \(X\) 上的 \(\pi\)-系统，则
\[
\lambda(\mathcal P)=\sigma(\mathcal P).
\]
因而，任何包含 \(\mathcal P\) 的 \(\lambda\)-系统都包含 \(\sigma(\mathcal P)\)。
\end{theorem}
\begin{proof}
令 \(\mathcal D=\lambda(\mathcal P)\)。证明的关键是把 \(\mathcal P\) 的交封闭性分两次推广到 \(\mathcal D\)。

先固定 \(P\in\mathcal P\)，并令
\[
\mathcal D_P
=\{\,A\in\mathcal D:A\cap P\in\mathcal D\,\}.
\]
下面逐条验证 \(\mathcal D_P\) 是 \(\lambda\)-系统。由于 \(X\cap P=P\in\mathcal D\)，有 \(X\in\mathcal D_P\)。若 \(A,B\in\mathcal D_P\) 且 \(A\subseteq B\)，则 \(B\setminus A\in\mathcal D\)，并且
\[
(B\setminus A)\cap P
=(B\cap P)\setminus(A\cap P)\in\mathcal D,
\]
故 \(B\setminus A\in\mathcal D_P\)。若 \(A_n\in\mathcal D_P\) 两两不交，则 \(A_n\cap P\) 也两两不交，而且
\[
\left(\bigcup_{n=1}^{\infty}A_n\right)\cap P
=\bigcup_{n=1}^{\infty}(A_n\cap P)\in\mathcal D.
\]
同时，\(\bigcup_nA_n\in\mathcal D\)，所以不交可数并仍属于 \(\mathcal D_P\)。

对每个 \(Q\in\mathcal P\)，\(\pi\)-系统的交封闭性给出 \(Q\cap P\in\mathcal P\subseteq\mathcal D\)，从而 \(Q\in\mathcal D_P\)。因此 \(\mathcal D_P\) 是包含 \(\mathcal P\) 的 \(\lambda\)-系统。由 \(\mathcal D=\lambda(\mathcal P)\) 的最小性，\(\mathcal D\subseteq\mathcal D_P\)；定义又给出反向包含，故 \(\mathcal D_P=\mathcal D\)。于是
\[
A\cap P\in\mathcal D
\quad(A\in\mathcal D,\ P\in\mathcal P).
\]

再固定 \(A\in\mathcal D\)，令
\[
\mathcal D_A
=\{\,B\in\mathcal D:A\cap B\in\mathcal D\,\}.
\]
这里 \(A\cap X=A\in\mathcal D\)，故 \(X\in\mathcal D_A\)。若 \(B,C\in\mathcal D_A\) 且 \(B\subseteq C\)，则 \(C\setminus B\in\mathcal D\)，并且
\[
A\cap(C\setminus B)
=(A\cap C)\setminus(A\cap B),
\]
右端属于 \(\mathcal D\)，所以 \(C\setminus B\in\mathcal D_A\)。若 \(B_n\in\mathcal D_A\) 两两不交，则 \(\bigcup_nB_n\in\mathcal D\)，且
\[
A\cap\bigcup_{n=1}^{\infty}B_n
=\bigcup_{n=1}^{\infty}(A\cap B_n)\in\mathcal D.
\]
所以 \(\mathcal D_A\) 是 \(\lambda\)-系统。上一段已经表明每个 \(P\in\mathcal P\) 都属于 \(\mathcal D_A\)，最小性再次给出 \(\mathcal D_A=\mathcal D\)。因此 \(\mathcal D\) 对有限交封闭。

由\cref{b1:prop:lambda-basic-closure}，\(\mathcal D\) 对补集封闭；结合有限交，它也对有限并封闭。对任意 \(A_n\in\mathcal D\)，令
\[
B_1=A_1,
\quad
B_n=A_n\setminus\bigcup_{k=1}^{n-1}A_k\quad(n\geq2).
\]
有限并、补集与有限交的封闭性保证
\(
B_n=A_n\cap\left(\bigcup_{k<n}A_k\right)^{\mathrm c}\in\mathcal D
\)，而 \(B_n\) 两两不交且 \(\bigcup_nB_n=\bigcup_nA_n\)。由 \(\lambda\)-系统的第三条，任意可数并仍在 \(\mathcal D\) 中。于是 \(\mathcal D\) 是包含 \(\mathcal P\) 的 \(\sigma\)-代数，故
\(
\sigma(\mathcal P)\subseteq\mathcal D
\)。另一方面，\(\sigma(\mathcal P)\) 本身是包含 \(\mathcal P\) 的 \(\lambda\)-系统，所以 \(\mathcal D\subseteq\sigma(\mathcal P)\)。两边相等。
\end{proof}

证明中两次固定集合并非重复：第一次把“与生成元相交”推广到 \(\mathcal D\)，第二次再把另一个变量也从生成元推广到 \(\mathcal D\)。这类“双重推广”会在乘积空间和测度唯一性中多次出现。

下面给出一个完全不使用测度概念的应用，它已经体现了生成元判别的力量。

\begin{corollary}
\label{b1:cor:inverse-images-agree}
设 \(f,g:Z\to Y\) 是两个映射，\(\mathcal P\) 是 \(Y\) 上的 \(\pi\)-系统。若
\[
f^{-1}(P)=g^{-1}(P)
\quad(P\in\mathcal P),
\]
则该等式对每个 \(B\in\sigma(\mathcal P)\) 都成立。
\end{corollary}
\begin{proof}
令
\[
\mathcal D
=\{\,B\subseteq Y:f^{-1}(B)=g^{-1}(B)\,\}.
\]
两个映射对 \(Y\) 的反像都是 \(Z\)，故 \(Y\in\mathcal D\)。若 \(A\subseteq B\) 且二者属于 \(\mathcal D\)，则反像与差运算交换，从而
\[
f^{-1}(B\setminus A)
=f^{-1}(B)\setminus f^{-1}(A)
=g^{-1}(B)\setminus g^{-1}(A)
=g^{-1}(B\setminus A).
\]
对两两不交集合的可数并，反像与并交换，同样保留等式。因此 \(\mathcal D\) 是包含 \(\mathcal P\) 的 \(\lambda\)-系统，结论由\cref{b1:thm:pi-lambda}得到。
\end{proof}

本章没有提前给出“概率测度的唯一性”作为推论，因为测度及其可数可加性要到下一章才正式定义。届时只需把“两个测度取值相等的集合”验证为 \(\lambda\)-系统，就能直接调用本节定理。

\subsection{单调类定理}
\label{b1:subsec:010304}

\(\lambda\)-系统仍然要求处理两两不交的可数并。有些极限论证天然只对递增列或递减列成立，这就引出更弱的单调闭性。

\begin{definition}
\label{b1:def:monotone-class}
集合族 \(\mathcal M\subseteq\mathcal P(X)\) 若满足：当 \(A_n\uparrow A\) 或 \(A_n\downarrow A\)，且每个 \(A_n\in\mathcal M\) 时，均有 \(A\in\mathcal M\)，则称 \(\mathcal M\) 为\term[dantiaolei]{单调类}[monotone class]。
\end{definition}

这里 \(A_n\uparrow A\) 表示 \(A_1\subseteq A_2\subseteq\cdots\) 且 \(\bigcup_nA_n=A\)；记号 \(A_n\downarrow A\) 表示反向包含且 \(\bigcap_nA_n=A\)。所有包含集合族 \(\mathcal E\) 的单调类之交仍是单调类；把这个最小单调类记为 \(\mathcal M(\mathcal E)\)。

单调类本身可能非常弱。例如在有限集合 \(X=\{1,2\}\) 上，集合族 \(\{\{1\}\}\) 已是单调类，因为其中的单调集合列只能是常值列；但它既不含 \(X\)，也不对补封闭。因此，下一定理要求起始生成族是集合代数，而不能是任意集合族。

\begin{theorem}
\label{b1:thm:monotone-class}
若 \(\mathcal A\) 是 \(X\) 上的集合代数，则
\[
\mathcal M(\mathcal A)=\sigma(\mathcal A).
\]
\end{theorem}
\begin{proof}
记 \(\mathcal M=\mathcal M(\mathcal A)\)。第一步证明 \(\mathcal M\) 对补集封闭。令
\[
\mathcal C
=\{\,E\in\mathcal M:E^{\mathrm c}\in\mathcal M\,\}.
\]
若 \(E_n\uparrow E\) 且每个 \(E_n\in\mathcal C\)，则 \(E_n^{\mathrm c}\downarrow E^{\mathrm c}\)；由 \(\mathcal M\) 的单调闭性，\(E^{\mathrm c}\in\mathcal M\)，故 \(E\in\mathcal C\)。递减列的情形相同，只是补集列递增。因此 \(\mathcal C\) 是单调类。集合代数 \(\mathcal A\) 对补封闭，所以 \(\mathcal A\subseteq\mathcal C\)。由 \(\mathcal M\) 的最小性，\(\mathcal M\subseteq\mathcal C\)，即 \(\mathcal C=\mathcal M\)。

第二步证明有限交封闭。先固定 \(A\in\mathcal A\)，令
\[
\mathcal M_A
=\{\,E\in\mathcal M:A\cap E\in\mathcal M\,\}.
\]
交运算与递增并、递减交交换，所以 \(\mathcal M_A\) 是单调类。对任意 \(B\in\mathcal A\)，集合代数的交封闭性给出 \(A\cap B\in\mathcal A\subseteq\mathcal M\)，故 \(\mathcal A\subseteq\mathcal M_A\)。最小性推出 \(\mathcal M_A=\mathcal M\)。于是
\[
A\cap E\in\mathcal M
\quad(A\in\mathcal A,\ E\in\mathcal M).
\]
现在固定 \(E\in\mathcal M\)，再令
\[
\mathcal N_E
=\{\,F\in\mathcal M:F\cap E\in\mathcal M\,\}.
\]
同样地，\(\mathcal N_E\) 是单调类；刚得到的结论说明 \(\mathcal A\subseteq\mathcal N_E\)，故 \(\mathcal N_E=\mathcal M\)。所以任意两个 \(\mathcal M\) 中的集合相交后仍属于 \(\mathcal M\)。

由补集和有限交的封闭性，\(\mathcal M\) 对有限并封闭。若 \(E_1,E_2,\ldots\in\mathcal M\)，则有限并
\(
F_n=\bigcup_{k=1}^{n}E_k
\)
属于 \(\mathcal M\)，且 \(F_n\uparrow\bigcup_kE_k\)。单调闭性给出 \(\bigcup_kE_k\in\mathcal M\)。于是 \(\mathcal M\) 是包含 \(\mathcal A\) 的 \(\sigma\)-代数，因而 \(\sigma(\mathcal A)\subseteq\mathcal M\)。反过来，\(\sigma(\mathcal A)\) 自身是包含 \(\mathcal A\) 的单调类；最小性给出 \(\mathcal M\subseteq\sigma(\mathcal A)\)。
\end{proof}

\subsection{函数型单调类定理}
\label{b1:subsec:010305}

集合型定理可以转化为函数上的推广工具。本小节的陈述会先使用下一节才正式命名的 Borel 可测性；第一次阅读时可以先掌握结论，读完第\ref{b1:subsec:010403}小节后再返回核对证明。为避免预先使用后章才系统发展的简单函数理论，下面把所需逼近直接写成有限个示性函数之和。

\begin{theorem}
\label{b1:thm:functional-monotone-class}
设 \(\mathcal A\) 是 \(X\) 上的集合代数，\(\mathcal H\) 是一族有界实值函数，满足：
\begin{enumerate}
  \item \label{b1:item:functional-monotone-1}
  \(\mathcal H\) 是实向量空间；
  \item \label{b1:item:functional-monotone-2}
  对每个 \(A\in\mathcal A\)，示性函数 \(\mathbf 1_A\) 属于 \(\mathcal H\)；
  \item \label{b1:item:functional-monotone-3}
  若 \(0\leq h_n\uparrow h\)，每个 \(h_n\in\mathcal H\)，且 \(h\) 有界，则 \(h\in\mathcal H\)。
\end{enumerate}
那么，每个满足
\[
f^{-1}(B)\in\sigma(\mathcal A)
\quad\bigl(B\in\mathcal B(\mathbb R)\bigr)
\]
的有界实值函数 \(f\) 都属于 \(\mathcal H\)。
\end{theorem}
\begin{proof}
先从函数类恢复一个集合类：
\[
\mathcal D
=\{\,E\subseteq X:\mathbf 1_E\in\mathcal H\,\}.
\]
假设 \(E_n\uparrow E\) 且每个 \(E_n\in\mathcal D\)。则
\(
0\leq\mathbf 1_{E_n}\uparrow\mathbf 1_E
\)，由第三条得到 \(E\in\mathcal D\)。又因 \(X\in\mathcal A\)，第二条给出常值函数 \(\mathbf 1_X\in\mathcal H\)。所以 \(E\in\mathcal D\) 时
\(
\mathbf 1_{E^{\mathrm c}}=\mathbf 1_X-\mathbf 1_E\in\mathcal H
\)，即 \(\mathcal D\) 对补集封闭。若 \(E_n\downarrow E\)，则 \(E_n^{\mathrm c}\uparrow E^{\mathrm c}\)；先对补集列使用递增情形，再取补，得到 \(E\in\mathcal D\)。因此 \(\mathcal D\) 是包含 \(\mathcal A\) 的单调类。由\cref{b1:thm:monotone-class}，
\[
\sigma(\mathcal A)\subseteq\mathcal D.
\]
换言之，每个 \(\sigma(\mathcal A)\) 中集合的示性函数都属于 \(\mathcal H\)。

现设 \(f\) 非负、有界，并满足定理中的反像条件。选取正整数 \(M\)，使 \(f(x)<M\) 对所有 \(x\in X\) 成立。定义
\[
s_n
=2^{-n}\sum_{k=1}^{M2^n}
\mathbf 1_{\{\,f\geq k2^{-n}\,\}}.
\]
闭半直线 \([k2^{-n},\infty)\) 是 Borel 集，所以求和中的每个集合都属于 \(\sigma(\mathcal A)\)；上一段与向量空间性质给出 \(s_n\in\mathcal H\)。对每个 \(x\in X\)，
\[
s_n(x)=2^{-n}\left\lfloor2^nf(x)\right\rfloor,
\]
故 \(0\leq s_n\uparrow f\)，且所有 \(s_n\) 都由 \(M\) 一致控制。第三条于是给出 \(f\in\mathcal H\)。

最后令 \(f\) 为任意满足反像条件的有界实值函数。选取常数 \(c\)，使 \(f+c\geq0\)。对任意 Borel 集 \(B\)，有
\[
(f+c)^{-1}(B)=f^{-1}(B-c),
\]
平移 \(x\mapsto x-c\) 是 \(\mathbb R\) 的同胚，故 \(B-c\) 仍是 Borel 集。因此 \(f+c\) 满足同一反像条件。由非负情形，\(f+c\in\mathcal H\)。又因 \(c\mathbf 1_X\in\mathcal H\)，向量空间性质给出
\(
f=(f+c)-c\mathbf 1_X\in\mathcal H
\)。
\end{proof}

定理中的反像条件将在下一节被正式称为 Borel 可测性。它的典型用法是：先对集合代数中的示性函数验证一个等式，再检查等式对线性组合和有界单调极限保持，便能把等式推广到所有有界可测函数。该证明只使用集合论闭性，不依赖积分理论。
